\section{LIT-SPECIALISTS: Behavioral Economics Specialists: Cross-Cutting Research}
\label{app:litad}

This appendix integrates papers from specialized behavioral economists
across multiple domains: Statman (behavioral finance), Rabin (game theory), Babcock
(labor economics), DellaVigna (field methods), Shafir (bounded rationality), Koeszegi
(reference dependence), Bowles (evolution), Haley (cross-cultural), Hogarth (judgment),
and Fox (ambiguity).

\subsection{Research Program Overview}

Behavioral Economics Specialists's research program addresses core behavioral economics themes:

\begin{itemize}
  \item \textbf{Behavioral Finance}: Investor psychology and market behavior
  \item \textbf{Game Theory}: Fairness and strategic interaction
  \item \textbf{Labor Markets}: Negotiation, gender, overconfidence
  \item \textbf{Bounded Rationality}: Cognitive limits and decision-making
  \item \textbf{Reference Dependence}: How expectations shape behavior
  \item \textbf{Cross-Cultural}: Variation in economic behavior across cultures
  \item \textbf{Judgment & Intuition}: Pattern-based and educated decisions
  \item \textbf{Ambiguity}: Uncertainty about probabilities
\end{itemize}

\subsection{Papers Integrated: 23 works}

\subsubsection{1993: Money Illusion...}

\paragraph{Citation.}
Shafir, Eldar, Diamond, Peter, et al. (1993). ``Money Illusion.''
\textit{Quarterly Journal of Economics}.

\paragraph{Core Finding.}
People fail to adjust for inflation; money illusion affects decisions

\paragraph{9C Integration.}
\begin{itemize}[nosep]
  \item \textbf{Domain}: Financial Behavior
  \item \textbf{Dimension}: F
  \item \textbf{Context}: Inflation Context
  \item \textbf{Complementarity}: γ = 0.71
\end{itemize}

\subsubsection{1995: Explaining Bargaining Impasse: The Role of Self-Se...}

\paragraph{Citation.}
Babcock, Linda, Loewenstein, George (1995). ``Explaining Bargaining Impasse: The Role of Self-Serving Biases.''
\textit{Journal of Economic Perspectives}.

\paragraph{Core Finding.}
Self-serving bias leads to breakdown of wage negotiations

\paragraph{9C Integration.}
\begin{itemize}[nosep]
  \item \textbf{Domain}: Labor Negotiations
  \item \textbf{Dimension}: S
  \item \textbf{Context}: Bargaining Context
  \item \textbf{Complementarity}: γ = 0.68
\end{itemize}

\subsubsection{1997: Building a Better Model of Bounded Rationality...}

\paragraph{Citation.}
Shafir, Eldar (1997). ``Building a Better Model of Bounded Rationality.''
\textit{SIAM Review}.

\paragraph{Core Finding.}
Bounded rationality models better predict behavior than rational choice

\paragraph{9C Integration.}
\begin{itemize}[nosep]
  \item \textbf{Domain}: Bounded Rationality
  \item \textbf{Dimension}: E
  \item \textbf{Context}: Cognitive Limits
  \item \textbf{Complementarity}: γ = 0.70
\end{itemize}

\subsubsection{1998: Psychology and Economics...}

\paragraph{Citation.}
Rabin, Matthew (1998). ``Psychology and Economics.''
\textit{Journal of Economic Literature}.

\paragraph{Core Finding.}
Behavioral economics integrates psychology into economic models

\paragraph{9C Integration.}
\begin{itemize}[nosep]
  \item \textbf{Domain}: Behavioral Economics
  \item \textbf{Dimension}: E
  \item \textbf{Context}: Psychological Context
  \item \textbf{Complementarity}: γ = 0.76
\end{itemize}

\subsubsection{2000: Risk Aversion and Expected-Utility Theory: A Calib...}

\paragraph{Citation.}
Rabin, Matthew (2000). ``Risk Aversion and Expected-Utility Theory: A Calibration Theorem.''
\textit{Econometrica}.

\paragraph{Core Finding.}
Standard risk aversion models inconsistent with observed behavior

\paragraph{9C Integration.}
\begin{itemize}[nosep]
  \item \textbf{Domain}: Risk Preferences
  \item \textbf{Dimension}: F
  \item \textbf{Context}: Uncertainty
  \item \textbf{Complementarity}: γ = 0.69
\end{itemize}

\subsubsection{2002: A Belief-Based Account of Decision Under Uncertain...}

\paragraph{Citation.}
Fox, Craig R., Tversky, Amos (2002). ``A Belief-Based Account of Decision Under Uncertainty.''
\textit{Management Science}.

\paragraph{Core Finding.}
Belief uncertainty (ambiguity) causes aversion distinct from risk aversion

\paragraph{9C Integration.}
\begin{itemize}[nosep]
  \item \textbf{Domain}: Ambiguity
  \item \textbf{Dimension}: E
  \item \textbf{Context}: Uncertainty
  \item \textbf{Complementarity}: γ = 0.71
\end{itemize}

\subsubsection{2002: Decisions Under Uncertainty: Empirical Studies of ...}

\paragraph{Citation.}
Shafir, Eldar (2002). ``Decisions Under Uncertainty: Empirical Studies of Behavior Under Risk and Ambiguity.''
\textit{Journal of Economic Literature}.

\paragraph{Core Finding.}
People show consistent patterns of deviation under uncertainty

\paragraph{9C Integration.}
\begin{itemize}[nosep]
  \item \textbf{Domain}: Decision Making
  \item \textbf{Dimension}: E
  \item \textbf{Context}: Uncertainty
  \item \textbf{Complementarity}: γ = 0.69
\end{itemize}

\subsubsection{2003: Biased Judgments of Fairness in Bargaining...}

\paragraph{Citation.}
Babcock, Linda (2003). ``Biased Judgments of Fairness in Bargaining.''
\textit{American Economic Review}.

\paragraph{Core Finding.}
Both sides see their own proposals as more fair than reality

\paragraph{9C Integration.}
\begin{itemize}[nosep]
  \item \textbf{Domain}: Labor Negotiations
  \item \textbf{Dimension}: S
  \item \textbf{Context}: Fairness Perception
  \item \textbf{Complementarity}: γ = 0.65
\end{itemize}

\subsubsection{2005: Educating Intuition...}

\paragraph{Citation.}
Hogarth, Robin M. (2005). ``Educating Intuition.''
\textit{Journal of Economic Literature}.

\paragraph{Core Finding.}
Intuition can be trained and improved through proper feedback loops

\paragraph{9C Integration.}
\begin{itemize}[nosep]
  \item \textbf{Domain}: Learning
  \item \textbf{Dimension}: E
  \item \textbf{Context}: Training Context
  \item \textbf{Complementarity}: γ = 0.64
\end{itemize}

\subsubsection{2006: A Model of Reference-Dependent Preferences...}

\paragraph{Citation.}
Koeszegi, Botond, Rabin, Matthew (2006). ``A Model of Reference-Dependent Preferences.''
\textit{Quarterly Journal of Economics}.

\paragraph{Core Finding.}
Formal model of reference dependence with testable predictions

\paragraph{9C Integration.}
\begin{itemize}[nosep]
  \item \textbf{Domain}: Reference Dependence
  \item \textbf{Dimension}: E
  \item \textbf{Context}: Expectation Context
  \item \textbf{Complementarity}: γ = 0.79
\end{itemize}

\subsubsection{2006: Behavioral Portfolio Theory...}

\paragraph{Citation.}
Statman, Meir (2006). ``Behavioral Portfolio Theory.''
\textit{Journal of Financial and Quantitative Analysis}.

\paragraph{Core Finding.}
Portfolio composition reflects psychological motives, not just financial optimization

\paragraph{9C Integration.}
\begin{itemize}[nosep]
  \item \textbf{Domain}: Portfolio Choice
  \item \textbf{Dimension}: F
  \item \textbf{Context}: Investment Context
  \item \textbf{Complementarity}: γ = 0.71
\end{itemize}

\subsubsection{2007: Prospect Theory and the Brain...}

\paragraph{Citation.}
Fox, Craig R., Poldrack, Russell A. (2007). ``Prospect Theory and the Brain.''
\textit{Neuroeconomics}.

\paragraph{Core Finding.}
Brain imaging reveals neural correlates of prospect theory operations

\paragraph{9C Integration.}
\begin{itemize}[nosep]
  \item \textbf{Domain}: Neuroeconomics
  \item \textbf{Dimension}: P
  \item \textbf{Context}: Neural Context
  \item \textbf{Complementarity}: γ = 0.68
\end{itemize}

\subsubsection{2008: Being Behaving: Biology, Behavior, and the Emergen...}

\paragraph{Citation.}
Bowles, Samuel (2008). ``Being Behaving: Biology, Behavior, and the Emergence of Social Preferences.''
\textit{Journal of Economic Literature}.

\paragraph{Core Finding.}
Social preferences evolved to facilitate cooperation in groups

\paragraph{9C Integration.}
\begin{itemize}[nosep]
  \item \textbf{Domain}: Social Preferences
  \item \textbf{Dimension}: S
  \item \textbf{Context}: Evolutionary Context
  \item \textbf{Complementarity}: γ = 0.73
\end{itemize}

\subsubsection{2009: Psychology and Economics: Evidence from the Field...}

\paragraph{Citation.}
DellaVigna, Stefano (2009). ``Psychology and Economics: Evidence from the Field.''
\textit{Journal of Economic Literature}.

\paragraph{Core Finding.}
Field evidence substantiates behavioral anomalies in real-world settings

\paragraph{9C Integration.}
\begin{itemize}[nosep]
  \item \textbf{Domain}: Methodology
  \item \textbf{Dimension}: E
  \item \textbf{Context}: Field Research
  \item \textbf{Complementarity}: γ = 0.74
\end{itemize}

\subsubsection{2009: Reference-Dependent Consumption Plans...}

\paragraph{Citation.}
Koeszegi, Botond, Rabin, Matthew (2009). ``Reference-Dependent Consumption Plans.''
\textit{American Economic Review}.

\paragraph{Core Finding.}
Reference points shape consumption and saving patterns

\paragraph{9C Integration.}
\begin{itemize}[nosep]
  \item \textbf{Domain}: Consumption Behavior
  \item \textbf{Dimension}: F
  \item \textbf{Context}: Expectation Effects
  \item \textbf{Complementarity}: γ = 0.76
\end{itemize}

\subsubsection{2010: Markets, Managers, and Messages: A Behavioral Fina...}

\paragraph{Citation.}
Statman, Meir (2010). ``Markets, Managers, and Messages: A Behavioral Finance Perspective.''
\textit{Financial Analysts Journal}.

\paragraph{Core Finding.}
Manager behavior and investor messaging shape market outcomes

\paragraph{9C Integration.}
\begin{itemize}[nosep]
  \item \textbf{Domain}: Behavioral Finance
  \item \textbf{Dimension}: F
  \item \textbf{Context}: Market Psychology
  \item \textbf{Complementarity}: γ = 0.65
\end{itemize}

\subsubsection{2011: The Causes of Gender Differences in Negotiation an...}

\paragraph{Citation.}
Babcock, Linda, Recalde, Maria P. (2011). ``The Causes of Gender Differences in Negotiation and Competing Explanations.''
\textit{Journal of Economic Literature}.

\paragraph{Core Finding.}
Gender differences in negotiation stem from social preferences and norms

\paragraph{9C Integration.}
\begin{itemize}[nosep]
  \item \textbf{Domain}: Gender Labor
  \item \textbf{Dimension}: S
  \item \textbf{Context}: Gender Norms
  \item \textbf{Complementarity}: γ = 0.72
\end{itemize}

\subsubsection{2011: The Yin and Yang of Cooperation in International J...}

\paragraph{Citation.}
Haley, Usha C.V., Fung, Chin-Meng (2011). ``The Yin and Yang of Cooperation in International Joint Ventures: A Culturally Contingent Model of Opportunism.''
\textit{Journal of International Business Studies}.

\paragraph{Core Finding.}
Cooperation behavior differs systematically across cultures

\paragraph{9C Integration.}
\begin{itemize}[nosep]
  \item \textbf{Domain}: Cross Cultural
  \item \textbf{Dimension}: S
  \item \textbf{Context}: Cultural Context
  \item \textbf{Complementarity}: γ = 0.65
\end{itemize}

\subsubsection{2011: Intuition: A Challenge to the Unbridged Mind...}

\paragraph{Citation.}
Hogarth, Robin M. (2011). ``Intuition: A Challenge to the Unbridged Mind.''
\textit{Yale University Press}.

\paragraph{Core Finding.}
Intuition is fast, pattern-based reasoning with reliable accuracy

\paragraph{9C Integration.}
\begin{itemize}[nosep]
  \item \textbf{Domain}: Intuition
  \item \textbf{Dimension}: E
  \item \textbf{Context}: Cognitive Processing
  \item \textbf{Complementarity}: γ = 0.67
\end{itemize}

\subsubsection{2012: The New Economics of Inequality and Redistribution...}

\paragraph{Citation.}
Bowles, Samuel (2012). ``The New Economics of Inequality and Redistribution.''
\textit{Cambridge University Press}.

\paragraph{Core Finding.}
Inequality and redistribution preferences are shaped by evolution and institutions

\paragraph{9C Integration.}
\begin{itemize}[nosep]
  \item \textbf{Domain}: Redistribution
  \item \textbf{Dimension}: S
  \item \textbf{Context}: Institutional Context
  \item \textbf{Complementarity}: γ = 0.68
\end{itemize}

\subsubsection{2013: An Integrative Theory of Green Organizational Inve...}

\paragraph{Citation.}
Haley, Usha C.V. (2013). ``An Integrative Theory of Green Organizational Investments: The Moderating Role of the Firm's Profitability and Institutional Context.''
\textit{Organization & Environment}.

\paragraph{Core Finding.}
Cultural and institutional factors shape sustainability decisions

\paragraph{9C Integration.}
\begin{itemize}[nosep]
  \item \textbf{Domain}: Sustainability
  \item \textbf{Dimension}: E
  \item \textbf{Context}: Institutional Context
  \item \textbf{Complementarity}: γ = 0.62
\end{itemize}

\subsubsection{2015: The Behavioral Economics of Taxation...}

\paragraph{Citation.}
DellaVigna, Stefano (2015). ``The Behavioral Economics of Taxation.''
\textit{Journal of Economic Literature}.

\paragraph{Core Finding.}
Behavioral factors significantly affect tax compliance and evasion

\paragraph{9C Integration.}
\begin{itemize}[nosep]
  \item \textbf{Domain}: Tax Behavior
  \item \textbf{Dimension}: F
  \item \textbf{Context}: Tax Context
  \item \textbf{Complementarity}: γ = 0.68
\end{itemize}

\subsubsection{2017: Finance for Normal People: How Investors and Marke...}

\paragraph{Citation.}
Statman, Meir (2017). ``Finance for Normal People: How Investors and Markets Behave.''
\textit{Oxford University Press}.

\paragraph{Core Finding.}
Investors seek happiness, not just returns; normal behavior differs from rational

\paragraph{9C Integration.}
\begin{itemize}[nosep]
  \item \textbf{Domain}: Behavioral Finance
  \item \textbf{Dimension}: F
  \item \textbf{Context}: Financial Context
  \item \textbf{Complementarity}: γ = 0.70
\end{itemize}

\subsection{Summary}

This appendix integrates 23 papers by Behavioral Economics Specialists
spanning research on behavioral finance, game theory, labor markets.
The papers demonstrate systematic patterns in human behavior that have important
implications for policy, organization, and individual decision-making.

\subsection{References}

For complete reference details and citations, see the master bibliography
\texttt{bcm2\_master\_references.bib}.

\vspace{1em}
\hrule
\vspace{0.5em}

\noindent\textit{Cross-references:}
Appendix G (Central Glossary), Chapter 14 (Applications)
